The Webster (major fractions) apportionment method rounds each state's
exact proportional quota to the nearest integer using the arithmetic
mean threshold of $s + 0.5$: a state receives $s+1$ seats if its
population-to-divisor ratio exceeds $s + 0.5$ seats, and $s$ seats
otherwise.
Webster is the principal academic benchmark for evaluating Huntington-Hill's
small-state bias, and it was used in U.S.\ congressional apportionment
in 1842 and from 1911 through 1940.
In proportional representation systems for party-list elections,
Webster's arithmetic mean rounding is known as the Sainte-Lagu\"e method,
the dominant PR formula in Scandinavia.

We prove that Webster minimises the sum of squared deviations
$\sum_i (s_i - q_i)^2$ from the exact quotas --- an optimality
property that makes it the divisor-method solution closest to
strict proportionality in the least-squares sense.
We show that Webster is paradox-immune as a divisor method and
that it is neutral between small and large states in the sense that
its arithmetic mean threshold neither systematically favors
small populations (as Huntington-Hill's geometric mean does slightly)
nor large populations (as Jefferson's floor does strongly).
Using 2020 Census data, Webster produces seat allocations identical
to Huntington-Hill for all 50 states;
we identify the state nearest the 0.5 rounding boundary
and quantify the population gap separating the two methods'
outcomes in 2020.
We document Webster's use in the 1842 Apportionment Act (the first
to mandate single-member districts) and its use under the 1911
Apportionment Act, and explain why Congress replaced it with
Huntington-Hill in 1941.
